\subsection{二次代数}

设 \(\alpha, \beta\) 是 A 的两个元素， \((e_1, e_2)\) 是 \(\mathbf{A}^2\) 的典范基。A 上 \((\alpha, \beta)\) 型的二次代数是指带有由乘法表（§ 1，第 7 号）定义的代数结构的 A-模 \(\mathbf{A}^2\)

\[
e _ {1} ^ {2} = e _ {1}, \qquad e _ {1} e _ {2} = e _ {2} e _ {1} = e _ {2}, \qquad e _ {2} ^ {2} = \alpha e _ {1} + \beta e _ {2}.\tag{2}
\]

一个与二次代数同构的A-代数E也称为二次代数。这等同于说E有一个由两个元素组成的基，其中一个为单位元。

可以证明，每个具有由两个元素组成的基的含单位元 A-代数都是二次代数（习题 1）。

若 A-代数的一个基 \((e_{1}, e_{2})\) 的乘法表为 (2)，则称它为 \((\alpha, \beta)\) 型的基。作为语言上的滥用，当一个二次代数具有 \((\alpha, \beta)\) 型的基时，就说它是 \((\alpha, \beta)\) 型的。

命题 1. 二次代数 E 是结合的且交换的。

E 的交换性由 (2) 中的关系式 \(e_{1}e_{2}=e_{2}e_{1}\) 得出；类似地，要验证结合性，只需验证 \(x(yz)=(xy)z\) 当 x, y, z 各取 \(e_{1}\) 或 \(e_{2}\) 时成立。现在，若 x, y, z 中至少有一个等于 \(e_{1}\) ，则此关系显然成立；对 x=y=z=e₂ 也成立，因为 E 交换；命题得证。

设 \(e\) 表示二次代数 \(\mathbf{E}\) 的单位元， \((e, i)\) 是 \(\mathbf{E}\) 的 \((\alpha, \beta)\) 型基； \(\mathbf{E}\) 的包含 \(e\) 的任何其他基因此形如 \((e, j)\) ，其中 \(j = \gamma e + \delta i\) （II，§ 7，第 2 号，命题 3 的推论）；此外， \((e, j)\) 为 \(\mathbf{E}\) 的基的必要且充分条件是 \(\delta\) 在 \(\mathbf{A}\) 中可逆；该条件显然充分；反之，若 \(i\) 是 \(i\) 在 \(\mathbf{E}/\mathbf{A}e\) 中的典范像，则 \(i\) 与 \(\bar{j} = \delta i\) 都必须构成 \(\mathbf{E}/\mathbf{A}e\) 的基，由此得该条件的必要性。于是

\[
j ^ {2} = (\gamma^ {2} + \alpha \delta^ {2}) e + (2 \gamma \delta + \beta \delta^ {2}) i = (\alpha \delta^ {2} - \gamma^ {2} - \beta \gamma \delta) e + (2 \gamma + \beta \delta) j;
\]

由此看出 E 是以下类型的

\[
(\alpha \delta^ {2} - \gamma^ {2} - \beta \gamma \delta , 2 \gamma + \beta \delta)\tag{3}
\]

对所有可逆的 \(\delta \in \mathbf{A}\) 以及所有 \(\gamma \in \mathbf{A}\) 。特别地，若 \(\mathbf{E}\) 是 \((\alpha, 2\beta')\) 型的，则 E 也是 \((\alpha + \beta'^2, 0)\) 型的，这可由取 \(\gamma = -\beta'\) 与 \(\delta = 1\) 看出。

命题 2. 设 E 是一个二次 A-代数，e 是它的单位元。对所有 \(u \in \mathbf{E}\) ，设 \(\mathrm{T}(u)\) 是自由 A-模 E 的自同态 \(m_u: x \mapsto ux\) 的迹（II，§ 4，第 3 号）。则映射 \(s\) ——定义为 \(s(u) = \mathrm{T}(u)\) . \(e - u\) ——是代数 E 的一个自同构，且 \(s^2(u) = u\) 对所有 \(u \in \mathbf{E}\) 成立。

设 \((e, i)\) 是 E 的 \((\alpha, \beta)\) 型基；则 \(T(e) = 2\) ，从而 \(s(e) = e\) ，且 \(T(i) = \beta\) ，从而 \(s(i) = \beta e - i\) 。于是 \((e, s(i))\) 是 E 的基，其类型由 (3) 给出，其中 \(\gamma = \beta\) 且 \(\delta = -1\) ，这重新给出 \((\alpha, \beta)\) ；由此得出 \(s\) 是代数 E 的自同构。由于 \(m_{s(u)} = sm_u s^{-1}\) ，A-模 E 的自同态 \(m_u\) 与 \(m_{s(u)}\) 具有相同的迹（II，§ 4，第 3 号，命题 3），因此

\[
s ^ {2} (u) = \mathrm{T} (u). e - s (u) = \mathrm{T} (u). e - (\mathrm{T} (u). e - u) = u
\]

对所有 \(u \in E\) 成立。

自同构 \(s\) 称为 A-代数 \(\mathbf{E}\) 的共轭，而 \(s(u)\) 称为 \(u\) 的共轭元。

若 \(u = \xi e + \eta i\) ，其中 \(\xi, \eta\) 属于 \(A\) ，则 \(s(u) = (\xi + \beta \eta)e - \eta i\) ，从而

\begin{enumerate}
\def\labelenumi{(\arabic{enumi})}
\setcounter{enumi}{3}
\tightlist
\item
\end{enumerate}

\[
\mathrm{T} (u) e = u + s (u) = (2 \xi + \beta \eta) e\tag{5}
\]

\[
u. s (u) = (\xi^ {2} + \beta \xi \eta - \alpha \eta^ {2}) e = N (u) e
\]

其中我们记 \(\mathrm{N}(u)=\xi^{2}+\beta\xi\eta-\alpha\eta^{2}\) 。元素 \(\mathrm{T}(u)\) 与 \(\mathrm{N}(u)\) （或当 A 与 Ae 典范地等同时的元素 \(\mathrm{T}(u)e\) 与 \(\mathrm{N}(u)e\) ）分别称为 u 的迹与范数。

当 \(\beta = 0\) 时，上述公式简化为

\[
s (\xi e + \eta i) = \xi e - \eta i, \quad \mathrm{T} (\xi e + \eta i) = 2 \xi , \quad \mathrm{N} (\xi e + \eta i) = \xi^ {2} - \alpha \eta^ {2}. \tag {6}
\]

显然，T 是 E 上的线性形式，N 是 E 上的二次形式（IX，§ 3，第 4 号）。由于 E 交换且结合，由 (5) 可得

\[
\mathrm{N} (u v) = \mathrm{N} (u)) \mathrm{N} (v).\tag{7}
\]

要使 u 在 E 中可逆，必要且充分的条件是 \(\mathrm{N}(u)\) 在 A 中可逆。因为由 \(\mathrm{N}(e)=1\) ，取 \(v=u^{-1}\) 便由 (7) 得出条件的必要性。反之，若 \(\mathrm{N}(u)\) 在 A 中可逆，则由 (5) 知 u 可逆，且

\[
u ^ {- 1} = (\mathrm{N} (u)) ^ {- 1} s (u).\tag{8}
\]

可以证明 \(\mathbf{N}(u)\) 是自同态 \(m_u\) 的行列式（§ 8，第 1 号）（参见 § 9 第 3 号，例 1）。

如下命题给出了交换域上二次代数的结构：

命题 3. 设 E 是一个 \((\alpha, \beta)\) 型的二次 A-代数。

\begin{enumerate}
\def\labelenumi{(\roman{enumi})}
\item
  若 \(\mathbf{A}\) 是一个域，且不包含任何元素 \(\zeta\) 使得 \(\zeta^2 = \alpha + \beta \zeta\) ，则 \(\mathbf{E}\) 是 \(a\) （交换）域（参见 V，§ 3）。
\item
  若环 A 含有元素 \(\zeta\) 使得 \(\zeta^2 = \alpha + \beta\zeta\) ，且 \(\beta - 2\zeta\) 可逆（相应地，为零），则 E 同构于 \(\mathbf{A} \times \mathbf{A}\) （相应地，是 \((0,0)\) 型的）。
\end{enumerate}

我们证明 (i)。设 \(\xi, \eta\) 是 A 的两个元素，且 \(u = \xi e + \eta i\) 。若 \(\eta \neq 0\) ，写 \(\theta = -\xi \eta^{-1}\) ，则由 (5) 得 \(\mathrm{N}(u) = \eta^2 (\theta^2 - \beta \theta - \alpha)\) ，于是由对 A 的假设得 \(\mathrm{N}(u) \neq 0\) ；若 \(\eta = 0\) ，则 \(\mathrm{N}(u) = \xi^2\) 。无论哪种情形，若 \(u \neq 0\) ，则 \(\mathrm{N}(u) \neq 0\) ，从而 \(\mathrm{N}(u)\) 在 A 中可逆，因而 \(u\) 在 E 中可逆。

我们现在证明 (ii)。典范基 \((e_{1}, e_{2})\) （代数 \(A \times A\) 的）是 \((0, 1)\) 型的。我们已见（公式 (3)）E 的类型为

\[
(\alpha \delta^ {2} - \gamma^ {2} - \beta \gamma \delta , 2 \gamma + \beta \delta)
\]

对所有 \(\gamma \in \mathbf{A}\) 以及所有 \(\delta\) （在 \(\mathbf{A}\) 中可逆的）。若 \(\beta - 2\zeta\) 可逆，取 \(\delta = (\beta - 2\zeta)^{-1}\) 与 \(\gamma = -\zeta(\beta - 2\zeta)^{-1}\) ；则 \(2\gamma + \beta\delta = 1\) ，且

\[
\alpha \delta^ {2} - \gamma^ {2} - \beta \gamma \delta = \delta^ {2} (\alpha - \zeta^ {2} + \beta \zeta) = 0;
\]

于是 E 是 \((0,1)\) 型的，从而同构于 A × A。若 \(\beta - 2\zeta = 0\) ，则已指出 E 是 \((\alpha + \zeta^{2}, 0)\) 型的，从而是 \((0,0)\) 型的，因为 \(\alpha + \zeta^{2} = 2\zeta^{2} - \beta\zeta = 0\) 。

\((0,0)\) 型的二次 A-代数也称为 A 上的对偶数代数。

